% -----------------------------------------------------------------------------------
\section{Commitment, evidence, and the within-problem map}\label{sec:commit}

\subsection{The commitment extension}

Reflective stability asks what a theory would do if, before any prediction is made, it could choose
which policy it will run. To state this as a criterion, the choice of a policy must itself be a
decision problem.

\begin{definition}[Commitment extension]\label{def:commit}
Let $D$ be a problem and $\lambda$ a \emph{coupling} of the laws $(\mu_q)_{q\in\Pol_D}$, that is, a
probability distribution on $\prod_{q\in\Pol_D}\supp\mu_q$ whose marginal on the $q$-th coordinate is
$\mu_q$. The \emph{commitment extension} $K_\lambda(D)$ is the problem in which
\begin{enumerate}[label=(\roman*)]
  \item the agent, at a single information set, chooses a policy $q\in\Pol_D$;
  \item this choice irrevocably installs $q$;
  \item the environment draws $\omega=(s_q)_{q\in\Pol_D}$ from $\lambda$, independently of the choice, and
        the agent receives $U_{K}(\omega,q)=U_D(s_q,q)$.
\end{enumerate}
Unless stated otherwise the structural states of $K_\lambda(D)$ are
$S_{K}=\prod_{q}\supp\mu_q$, and $K(D)$ denotes the product coupling
$\lambda=\bigotimes_q\mu_q$. The constructor is tagged and well founded, so that $K(D)$,
$K(K(D))$, \dots{} have unambiguous decompositions, and $\cls^*$ is closed under it.
\end{definition}

The extension models a \emph{causal} choice of policy made before prediction: changing the
commitment changes the environment's subsequent response, because the coordinate $s_q$ that is
paid out depends on $q$, while the law $\lambda$ of the whole vector does not.

\begin{proposition}\label{prop:commitvalue}
For every coupling $\lambda$ and every $q\in\Pol_D$, $V_{K_\lambda(D)}(q)=V_D(q)$.
\end{proposition}

\begin{proof}
$V_{K_\lambda(D)}(q)=\E_\lambda\bigl[U_D(s_q,q)\bigr]=\E_{\mu_q}\bigl[U_D(s,q)\bigr]=V_D(q)$, since the
$q$-th marginal of $\lambda$ is $\mu_q$.
\end{proof}

\begin{definition}[Commitment stability, \cri{4}]
$T$ satisfies \cri{4} on $D$ if the commitment $q=T_{K(D)}$ that $T$ chooses in $K(D)$ satisfies
$q(h)=T_D(h)$ at every information set $h$ positively reached under $T_D$.
\end{definition}

Unlike the other criteria, \cri{4} relates a theory's behaviour on two different problems. It is
the first cross-problem condition in this paper, and the only one stated without passing to the
reduced form of Section~\ref{sec:matrix}.

\subsection{The map of \texorpdfstring{\cri{1}--\cri{5}}{C1-C5}}

The consistent combinations need explicit total witnesses, not merely examples on a few problems.

\begin{proposition}[A coherent ex-ante theory]\label{prop:coherent}
There is a decision theory on $\cls^*$ satisfying \cri{1} and \cri{4}.
\end{proposition}

\begin{proof}
On problems that are not commitment extensions, let $T_D=T^{\mathrm{ex}}_D$. Recursively, in $K(D)$
let $T$ commit to the policy $T_D$ already selected for $D$. By Proposition~\ref{prop:commitvalue},
$T_D$ maximises $V_{K(D)}$, so \cri{1} holds on $K(D)$; \cri{4} holds with exact agreement.
\end{proof}

\cri{1} alone does not imply \cri{4}: on a problem in which both actions yield zero, an ex-ante
maximising lookup table may choose different tied actions in $D$ and in $K(D)$.

\begin{proposition}[A dominance-respecting theory]\label{prop:nu}
For each $D$ fix a full-support distribution $\nu_D$ on $S_D$, for instance the uniform one, and let
\[
  T^{\nu}_D=\first\argmax_{\pi\in\Pol_D}\;\sum_{s\in S_D}\nu_D(s)\,U_D(s,\pi).
\]
Then $T^{\nu}$ satisfies \cri{2} on every problem.
\end{proposition}

\begin{proof}
Let $\pi=T^\nu_D$ and suppose that at an information set $h$ positively reached under $\pi$ some $a$
satisfies $U_D(s,\dev{\pi}{h}{a})>U_D(s,\pi)$ for all $s\in R_h(\pi)$. By Lemma~\ref{lem:local}(ii)
the payoff is unchanged on $S_D\setminus R_h(\pi)$. The set $R_h(\pi)$ is nonempty because
$r_\pi(h)>0$, and $\nu_D$ gives it positive weight, so $\dev{\pi}{h}{a}$ strictly improves the
objective that $\pi$ maximises, a contradiction.
\end{proof}

$T^\nu$ is an existence witness, not a recommended theory: its fixed weights ignore the
predictive dependence altogether.

\begin{proposition}\label{prop:c2c4}
No theory satisfies \cri{2} and \cri{4} on a class containing Newcomb's problem with $\varepsilon=0$
and its commitment extension.
\end{proposition}

\begin{proof}
In $D$, \cri{2} requires $B$ (Proposition~\ref{prop:newcomb}). At $\varepsilon=0$ every support is a
singleton, so $K(D)$ has the single structural state (opaque box full under $O$, empty under $B$),
and committing to $O$ yields $2$ while committing to $B$ yields $1$. Committing to $O$ therefore
strictly dominates, and \cri{2} in $K(D)$ requires it; \cri{4} requires the commitment to agree with
$B$ at the always-reached information set.
\end{proof}

\begin{theorem}[Within-problem map]\label{thm:map15}
On $\cls^*$ the inclusion-minimal inconsistent subsets of $\{\cri{1},\dots,\cri{5}\}$ are exactly
\[
  \{\cri{3}\},\qquad\{\cri{1},\cri{2}\},\qquad\{\cri{2},\cri{4}\},
\]
and the maximal consistent subsets are exactly $\{\cri{1},\cri{4},\cri{5}\}$, witnessed by
Proposition~\ref{prop:coherent}, and $\{\cri{2},\cri{5}\}$, witnessed by
Proposition~\ref{prop:nu}.
\end{theorem}

\begin{proof}
The three sets are inconsistent by Propositions~\ref{prop:did}, \ref{prop:newcomb}, and
\ref{prop:c2c4}; their proper subsets are consistent by the two witnesses (\cri{5} holds for every
theory). A consistent set cannot contain \cri{3}. If it contains \cri{2}, it contains neither \cri{1}
nor \cri{4}, so it lies in $\{\cri{2},\cri{5}\}$. Otherwise its members lie in
$\{\cri{1},\cri{4},\cri{5}\}$.
\end{proof}

Among \cri{1}--\cri{5}, then, there is no minimally inconsistent family of three or more
individually satisfiable criteria. The requested minimality is attained only by the pairs
$\{\cri{1},\cri{2}\}$ and $\{\cri{2},\cri{4}\}$, whose singletons have explicit total witnesses, and
the first is a direct pointwise disagreement in the sense of Lemma~\ref{lem:pointwise}.

\subsection{An evidential criterion}\label{sec:evidential}

An evidential criterion requires a joint law of acts and states. Under a deterministic policy,
conditioning on an alternative action is ordinarily conditioning on a null event, so the expression
$\E[\text{payoff}\mid h,\text{agent takes }a]$ cannot be read off the actual distribution of play.
The following completion makes it determinate.

\begin{definition}[Reference-policy evidential criterion, \cri{6}$^\kappa$]
For each $D$ fix a full-support distribution $\kappa_D$ on $\Pol_D$ and define the joint law
$J_D(\Pi=\pi,S=s)=\kappa_D(\pi)\mu_\pi(s)$, with $Y=U_D(S,\Pi)$. For $h\in H_D$ and $a\in A_h$ let
\[
  e_h(a)=\E_{J_D}\bigl[Y\mid h\text{ reached},\ \Pi(h)=a\bigr],
\]
with $e_h(a)=-\infty$ when the conditioning event is null. $T$ satisfies \cri{6}$^\kappa$ on $D$ if
$T_D(h)\in\argmax_a e_h(a)$ at every $h$ positively reached under $T_D$.
\end{definition}

This is a specified evidential criterion, not an identification with every version of evidential
decision theory; in sequential problems it evaluates continuation behaviour with the reference
distribution over policies. It is individually satisfiable: choose an $e_h$-maximiser at every $h$
with $J_D(h\text{ reached})>0$ and anything elsewhere. If $h$ is positively reached under $T_D$, then
$J_D(h\text{ reached})\ge\kappa_D(T_D)\,r_{T_D}(h)>0$, so the constraint is met, and some action has
a finite evidential value there.

\begin{proposition}\label{prop:c6}
\begin{enumerate}[label=(\roman*)]
  \item In Newcomb's problem $e(O)=V(O)$ and $e(B)=V(B)$, so \cri{6}$^\kappa$ requires $O$ for
        $\varepsilon<\tfrac14$ and conflicts with \cri{2}.
  \item In counterfactual mugging \cri{6}$^\kappa$ requires refusing, since $e(\text{Pay})=-c<0=e(\text{Refuse})$
        at the tails information set, and conflicts with \cri{1}.
  \item In the commitment extension of counterfactual mugging, $e(q)=V_D(q)$, so \cri{6}$^\kappa$
        commits to paying while refusing in the original problem, and conflicts with \cri{4}.
\end{enumerate}
\end{proposition}

\begin{proof}
(i) With one always-reached information set, conditioning on the action identifies the policy, so
$e(a)=\E_{\mu_a}U(\cdot,a)=V(a)$. (ii) Conditional on tails and on the tails prescription being
\enquote{pay}, the payoff is $-c$; on \enquote{refuse} it is $0$. (iii) In $K(D)$ the commitment
is the policy variable and $\lambda$ does not depend on it, so $e(q)=\E_\lambda U_D(s_q,q)=V_D(q)$, and
$V_D(\text{Pay})>V_D(\text{Refuse})$; the tails information set is reached with probability
$\tfrac12$ under the refusing policy.
\end{proof}

\begin{theorem}[Map with the evidential criterion]\label{thm:map16}
On $\cls^*$ the inclusion-minimal inconsistent subsets of $\{\cri{1},\dots,\cri{5},\cri{6}^\kappa\}$
are
\[
  \{\cri{3}\},\quad\{\cri{1},\cri{2}\},\quad\{\cri{2},\cri{4}\},\quad
  \{\cri{1},\cri{6}^\kappa\},\quad\{\cri{2},\cri{6}^\kappa\},\quad\{\cri{4},\cri{6}^\kappa\},
\]
and the maximal consistent subsets are $\{\cri{1},\cri{4},\cri{5}\}$, $\{\cri{2},\cri{5}\}$, and
$\{\cri{6}^\kappa,\cri{5}\}$.
\end{theorem}

\begin{proof}
Among \cri{1}, \cri{2}, \cri{4}, \cri{6}$^\kappa$, every pair except $\{\cri{1},\cri{4}\}$ is
inconsistent by Propositions~\ref{prop:newcomb}, \ref{prop:c2c4}, and~\ref{prop:c6}, and each listed
maximal set has a total witness.
\end{proof}

A different evidential semantics can change this table; Section~\ref{sec:evsem} compares two. The
theorem is a statement about the displayed definitions.

\subsection{Noisy commitment and couplings}\label{sec:couplings}

Proposition~\ref{prop:c2c4} uses exact prediction and, implicitly, a choice of coupling. Both
matter, and they matter differently for structural and for expectational comparisons. Let the
opaque prize be $m$ and the transparent bonus $c$, with $m>c>0$, and let the symmetric error be
$\varepsilon\in[0,1]$. Write $g=m(1-2\varepsilon)-c$ for the ex-ante advantage of $O$.

Let $X$ indicate a full opaque box under commitment to $O$, and $Y$ a full opaque box under commitment
to $B$, so that $\Pr(X=1)=1-\varepsilon$ and $\Pr(Y=1)=\varepsilon$. Every coupling is determined by
$t=\Pr(X=1,Y=1)$, which ranges over $0\le t\le\min(\varepsilon,1-\varepsilon)$:

\begin{center}
\small
\begin{tabular}{@{}ccccc@{}}
  \toprule
  $(X,Y)$ & probability & payoff of $O$ & payoff of $B$ & difference $O-B$ \\
  \midrule
  $(1,1)$ & $t$ & $m$ & $m+c$ & $-c$ \\
  $(1,0)$ & $1-\varepsilon-t$ & $m$ & $c$ & $m-c$ \\
  $(0,1)$ & $\varepsilon-t$ & $0$ & $m+c$ & $-m-c$ \\
  $(0,0)$ & $t$ & $0$ & $c$ & $-c$ \\
  \bottomrule
\end{tabular}
\end{center}

\begin{proposition}[Noisy commitment in Newcomb's problem]\label{prop:noisycommit}
\begin{enumerate}[label=(\roman*)]
  \item For every coupling, the expected causal value of committing to $q$ is $V_D(q)$ whatever the
        agent's beliefs about its own commitment; every criterion that ranks commitments by expected
        causal value requires commitment to $O$ exactly when $g>0$, while the same criterion requires
        $B$ in $D$.
  \item With structural states $S_K=\supp\lambda$, commitment to $O$ strictly dominates commitment to $B$
        if and only if $\varepsilon=0$; commitment to $B$ strictly dominates if and only if
        $\varepsilon\ge\tfrac12$ and $t=1-\varepsilon$.
  \item With structural states $S_K=\prod_q\supp\mu_q$, neither commitment dominates for
        $0<\varepsilon<1$, whatever the coupling.
\end{enumerate}
Consequently the \cri{2}--\cri{4} conflict of Proposition~\ref{prop:c2c4} occurs, on this pair of
problems, only at $\varepsilon=0$ and under every coupling, while the reversal of expectational
comparisons between $D$ and $K_\lambda(D)$ occurs whenever $g>0$, independently of $\lambda$.
\end{proposition}

\begin{proof}
(i) Since $\lambda$ does not depend on the commitment, the causal and the ex-ante value of committing
to $q$ coincide and equal $V_D(q)$ by Proposition~\ref{prop:commitvalue}. In $D$ the bonus $c$ is
gained in every state, so every expectational causal comparison favours $B$. (ii) $O$ dominates
only if the three rows with negative difference are null: $t=0$ and $\varepsilon-t=0$, so
$\varepsilon=0$; conversely, at $\varepsilon=0$ only the row $(1,0)$ remains. Directly: at positive error
$O$ pays $0$ with positive probability whereas $B$ always pays at least $c>0$, so no coupling makes
$O$ even weakly dominant. $B$ dominates exactly when the row $(1,0)$ is null, $t=1-\varepsilon$, which
the constraint $t\le\min(\varepsilon,1-\varepsilon)$ allows only for $\varepsilon\ge\tfrac12$. (iii) For
$0<\varepsilon<1$ all four rows are structural states, and the differences have both signs.
\end{proof}

Changing a probability law cannot change \cri{2} while its structural state set is held fixed; the
coupling matters for \cri{2} only through the convention $S_K=\supp\lambda$. The three couplings one
might call natural give
\begin{center}
\small
\begin{tabular}{@{}lc@{}}
  \toprule
  coupling & $t$ \\
  \midrule
  independent (product) & $\varepsilon(1-\varepsilon)$ \\
  shared predictor-correctness bit & $0$ \\
  comonotone realised payoffs & $\min(\varepsilon,1-\varepsilon)$ \\
  \bottomrule
\end{tabular}
\end{center}
and none of them restores dominance of $O$ at positive error. Couplings do matter in other problems:
for a shared Bernoulli variable $Z$, the payoffs $U=1+100Z$ and $W=100Z$ satisfy $U>W$ everywhere on
the support of the shared coupling, whereas the product coupling admits the reversing pair $U=1$,
$W=100$. Which couplings produce dominance is settled in general by the following criterion, a
finite form of Strassen's theorem on stochastic order \autocite{strassen1965}.

\begin{proposition}[Dominance under some coupling]\label{prop:strassen}
Let $\nu_U,\nu_W$ be finitely supported laws on $\Rl$ with distribution functions $F_U,F_W$.
\begin{enumerate}[label=(\roman*)]
  \item There is a coupling with $U\ge W$ almost surely if and only if $F_U(x)\le F_W(x)$ for all $x$.
  \item There is a coupling with $U>W$ almost surely if and only if $F_U(x)\le F_W(x^-)$ for every
        $x\in\supp\nu_U$.
\end{enumerate}
\end{proposition}

\begin{proof}
Necessity follows from the inclusions $\{U\le x\}\subseteq\{W\le x\}$ and
$\{U\le x\}\subseteq\{W<x\}$, which hold almost surely under the respective couplings. For
sufficiency, let $V$ be uniform on $(0,1]$ and set $U=F_U^{-1}(V)$, $W=F_W^{-1}(V)$, with
$F^{-1}(v)=\min\{y:F(y)\ge v\}$. In case (i), $F_U\le F_W$ gives $F_W^{-1}\le F_U^{-1}$ pointwise.
In case (ii), let $x=F_U^{-1}(v)\in\supp\nu_U$; then $F_W(x^-)\ge F_U(x)\ge v$, and since $\nu_W$ is
finitely supported, $F_W(y)\ge v$ for some $y<x$, so $F_W^{-1}(v)<x$.
\end{proof}

A payoff coupling lifts to a coupling of the underlying states by distributing mass within payoff
fibres according to the conditional state laws. For a general problem, the residual content of
\cri{4} is then an agreement condition: some \cri{2}-permitted original policy must agree, on the
information sets it positively reaches, with a commitment that no other commitment structurally
dominates. Proposition~\ref{prop:strassen} decides which commitments are dominated; it does not by
itself decide whether \cri{2} and \cri{4} are consistent on a whole class of problems with
imperfect prediction, which remains open (Section~\ref{sec:conclusion}).
